\subsubsection{从自对偶压缩到一变量 \texorpdfstring{$\Xi$}{Xi}：严格函数方程与临界线实值性}\label{subsec:xi-from-self-dual}
本节把附录中“自对偶 + 完成化 + 单位圆实值化”的三件套压缩为一个一变量 \(\Xi\)-对象，并把“临界线”解释为 \(u\)-参数的单位圆切片。

\paragraph{自对偶假设}
设 \(\Delta(z,u)\) 是二变量多项式/解析函数（在本文的核口径下典型为
\(\Delta(z,u)=\det(I-zB(u))\)，见 \S\ref{subsec:zeta-online-kernel}），并满足二变量自对偶恒等式
\[
\boxed{\ \Delta(z,u)=\Delta(uz,1/u)\ }.
\]
对同步核加权情形，该恒等式在附录 \ref{app:kernel-self-dual} 给出；并由此导出完成化偶性与系数回文约束（例如推论 \ref{cor:delta-coeff-palindrome}）。

\paragraph{新的一变量对象：任意底 \texorpdfstring{$L>1$}{L>1}}
取任意常数底 \(L>1\)，定义
\[
z_L(s):=L^{-s},\qquad u_L(s):=L^{2s-1},
\]
并定义一变量函数
\[
\boxed{\ \Xi_{\Delta,L}(s):=\Delta\!\bigl(z_L(s),u_L(s)\bigr)=\Delta\!\bigl(L^{-s},L^{2s-1}\bigr)\ }.
\]
当取 \(L=\lambda=\rho(B(1))\) 时，这与附录 \ref{app:kernel-functional-equation} 的截面 \(\Xi(s)\) 一致（命题 \ref{prop:kernel-xi-functional-equation}）。

\begin{proposition}[严格函数方程：自对偶 \(\Rightarrow\ s\mapsto 1-s\)]\label{prop:xi-delta-L-functional-equation}
在上述自对偶假设下，对任意 \(L>1\) 都有严格恒等式
\[
\boxed{\ \Xi_{\Delta,L}(s)=\Xi_{\Delta,L}(1-s)\ }.
\]
\end{proposition}
\begin{proof}
由自对偶恒等式，
\[
\Xi_{\Delta,L}(s)=\Delta(L^{-s},L^{2s-1})
=\Delta(L^{2s-1}L^{-s},L^{1-2s})
=\Delta(L^{s-1},L^{1-2s}).
\]
而 \(z_L(1-s)=L^{-(1-s)}=L^{s-1}\)、\(u_L(1-s)=L^{2(1-s)-1}=L^{1-2s}\)，故结论成立。
\end{proof}

\paragraph{完成化与“临界线实值性”}
写
\[
\Delta(z,u)=\sum_{j=0}^{N}c_j(u)\,z^j.
\]
若系数满足完成化回文结构 \(c_j(u)=u^j c_j(1/u)\)（推论 \ref{cor:delta-coeff-palindrome}），则可定义完成化系数
\[
d_j(u):=u^{-j/2}c_j(u).
\]
附录推论 \ref{cor:completion-real-line} 表明：当 \(|u|=1\) 时，\(d_j(u)\in\RR\)（完成化后“单位圆上系数为实”）。

\begin{proposition}[临界线实值性：完成化 \(\Rightarrow\) Hardy 型实值化]\label{prop:xi-delta-L-critical-line-real}
在自对偶假设并满足推论 \ref{cor:completion-real-line} 的条件下，对任意 \(L>1\) 与任意实数 \(t\) 有
\[
\boxed{\ \Xi_{\Delta,L}\!\left(\tfrac12+it\right)\in\RR\ }.
\]
\end{proposition}
\begin{proof}
取 \(s=\tfrac12+it\)，则 \(u_L(s)=L^{2it}\) 满足 \(|u_L(s)|=1\)，并且
\[
\Xi_{\Delta,L}(s)
=\Delta(L^{-s},u_L(s))
=\sum_{j=0}^{N}c_j(u_L(s))\,L^{-js}
=\sum_{j=0}^{N}d_j(u_L(s))\,\bigl(L^{-1/2}\bigr)^{j}\,e^{-ijt\log L}.
\]
由 \(|u_L(s)|=1\Rightarrow d_j(u_L(s))\in\RR\)，上述表达是一个关于实变量 \(e^{-it\log L}\) 的“实系数”三角多项式在 \(t\) 处的取值，因此其值为实数。更直接地，等价于把附录中 Hardy 实值化 \(Z(t)\) 的证明（命题 \ref{prop:kernel-hardy-real-even}）将 \(\log\lambda\) 替换为 \(\log L\)。
\end{proof}

\begin{remark}[临界线 \texorpdfstring{$\Leftrightarrow$}{<->} 单位圆切片]\label{rem:xi-critical-line-unit-circle}
由 \(u_L(s)=L^{2s-1}\) 可见：直线 \(\Re(s)=\tfrac12\) 恰对应 \(|u_L(s)|=1\)。
因此对该 \(\Xi_{\Delta,L}\) 而言，“临界线”并非额外结构，而是自对偶对称在参数侧的固定几何：单位圆切片。
\end{remark}

\paragraph{尺度化观点：最坏扭曲慢模的扩散尺度}
在有限模数 \(m\) 的扭曲（例如 \(u=\omega_m^r=e^{2\pi ir/m}\)）中，记最坏非平凡谱半径比为
\[
\eta_m:=\max_{1\le r\le m-1}\frac{\rho\bigl(B(\omega_m^r)\bigr)}{\rho(B(1))}\in(0,1].
\]
由于 \(\widehat{\ZZ}\) 的对偶包含任意高阶角色，必存在 \(m\to\infty\) 使 \(\omega_m\to 1\)，因此一般不可能在全部 \(m\) 上得到统一的谱隙下界 \(\eta_m\le 1-\delta\)（\(\delta>0\)）。
同步核加权情形中，附录已给出 near-coboundary 小角二次律及其扩散时间尺度：例如推论 \ref{cor:sync-rho-gap-m2} 与推论 \ref{cor:sync-rho-gap-m8-diffusive} 给出
\[
\eta_m\approx \exp\!\Bigl(-\kappa\,(2\pi/m)^2\Bigr),
\qquad
\eta_m^n\approx \exp\!\Bigl(-\kappa(2\pi)^2\,\frac{n}{m^2}\Bigr),
\]
从而“要在模 \(m\) 上看到稳定等分布”自然对应扩散阈值 \(n\gg m^2\)。

\begin{conjecture}[尺度化核版 GRH/Chebotarev（双极限扩散口径）]\label{conj:scaled-kernel-grh-diffusive}
存在常数 \(c>0\) 与 \(\kappa>0\)，使得对所有模数 \(m\ge 2\) 与所有长度 \(n\ge 1\)，任意由模 \(m\) 寄存器定义的同余 cylinder 事件的等分布误差满足
\[
\boxed{\ \mathrm{Err}(n,m)\ \le\ C(m)\,\exp\!\Bigl(-c\,\frac{n}{m^2}\Bigr)\ },
\]
其中 \(C(m)\) 至多多项式增长，并且主指数 \(c\) 与二次律常数 \(\kappa\) 同阶（在可解析的核族中可由压力曲率/方差密度给出）。
因此在区域 \(m\le n^{1/2-\varepsilon}\)（任意固定 \(\varepsilon>0\)）内得到一致的“平方根级可见主项”。
\end{conjecture}

